\documentclass[12pt,a4paper]{article}

\usepackage[spanish]{babel}
\usepackage[utf8]{inputenc}
\usepackage[T1]{fontenc}
\usepackage{geometry}
\usepackage{setspace}
\usepackage{graphicx}
\usepackage{tikz}
% CARGA ESTA LIBRERÍA PARA EVITAR EL ERROR
\usetikzlibrary{babel, positioning} 
\usepackage{enumitem}
\usepackage{array}

\geometry{margin=2.5cm}
\setstretch{1.5}

\title{Gestión del Conocimiento en Ingeniería de Sistemas}
\author{}
\date{}

\begin{document}

\maketitle

\section{Gestión del Conocimiento}

\subsection{Definición}
La gestión del conocimiento (Knowledge Management, KM) es un enfoque sistemático mediante el cual las organizaciones crean, capturan, organizan, comparten y utilizan el conocimiento con el propósito de mejorar su desempeño. Este proceso integra personas, tecnología y procesos, permitiendo aprovechar tanto el conocimiento individual como el colectivo.

Desde una perspectiva académica, puede definirse como un conjunto de prácticas orientadas a identificar, estructurar y difundir el conocimiento dentro de una organización para generar valor y ventaja competitiva.

\subsection{Importancia en las organizaciones}
La gestión del conocimiento es fundamental debido a su impacto directo en la eficiencia organizacional. Entre sus principales beneficios se destacan:

\begin{itemize}
    \item Previene la pérdida de conocimiento crítico.
    \item Mejora la toma de decisiones mediante información estructurada.
    \item Fomenta la innovación al facilitar la integración de ideas.
    \item Optimiza procesos, reduciendo errores y tiempos.
    \item Promueve el aprendizaje organizacional continuo.
\end{itemize}

En ingeniería de sistemas, estos beneficios se reflejan en mejores prácticas de desarrollo, documentación eficiente y reutilización de soluciones.

\subsection{Objetivo principal}
El objetivo central de la gestión del conocimiento es transformar el conocimiento en un activo estratégico reutilizable, compartido y aplicable para mejorar el rendimiento organizacional. Esto implica garantizar que la información sea accesible, comprensible y útil.

\section{Conocimiento Tácito vs Conocimiento Explícito}

\subsection{Conocimiento tácito}
El conocimiento tácito reside en la mente de las personas y se basa en experiencias, habilidades y percepciones. Es difícil de formalizar y transmitir.

\textbf{Ejemplo:} Un programador senior que depura errores complejos sin seguir un procedimiento documentado.

\subsection{Conocimiento explícito}
El conocimiento explícito puede ser documentado, estructurado y compartido fácilmente mediante diversos medios como manuales, bases de datos o diagramas.

\textbf{Ejemplo:} Documentación técnica de una API.

\subsection{Características principales}

\subsubsection{Conocimiento tácito}
\begin{itemize}
    \item Personal y subjetivo
    \item Difícil de formalizar
    \item Basado en la experiencia
    \item Transmisión mediante interacción directa
\end{itemize}

\subsubsection{Conocimiento explícito}
\begin{itemize}
    \item Formal y estructurado
    \item Fácil de almacenar
    \item Independiente del individuo
    \item Documentado en medios físicos o digitales
\end{itemize}

\subsection{Ejemplos comparativos}

\begin{center}
\begin{tabular}{|m{4cm}|m{4cm}|m{4cm}|}
\hline
\textbf{Contexto} & \textbf{Conocimiento tácito} & \textbf{Conocimiento explícito} \\
\hline
Software & Intuición para optimizar código & Estándares de programación \\
\hline
Medicina & Experiencia del cirujano & Protocolos médicos \\
\hline
Educación & Habilidad para enseñar & Plan de estudios \\
\hline
\end{tabular}
\end{center}

\subsection{Comparación conceptual}
La diferencia fundamental radica en su almacenamiento y transmisión:

\begin{itemize}
    \item El conocimiento tácito reside en las personas.
    \item El conocimiento explícito reside en los sistemas.
    \item El tácito es difícil de capturar pero altamente valioso.
    \item El explícito es fácil de compartir pero puede perder contexto.
\end{itemize}

\section{Transformación del Conocimiento}

\subsection{Descripción del proceso}
La transformación del conocimiento tácito en explícito implica convertir experiencias individuales en información estructurada. Este proceso es fundamental en modelos como el SECI.

\subsection{Etapas del proceso}

\begin{enumerate}
    \item Identificación del conocimiento crítico
    \item Extracción mediante entrevistas u observación
    \item Articulación en conceptos
    \item Codificación en documentos
    \item Validación
    \item Almacenamiento y difusión
\end{enumerate}

\subsection{Diagrama del proceso}

\begin{center}
\resizebox{\textwidth}{!}{%
\begin{tikzpicture}[node distance=1.5cm, every node/.style={draw, rectangle, rounded corners, align=center, font=\small}]
\node (id) {Identificación};
\node (ex) [right=of id] {Extracción};
\node (ar) [right=of ex] {Articulación};
\node (co) [right=of ar] {Codificación};
\node (va) [right=of co] {Validación};
\node (al) [right=of va] {Almacenamiento};

\draw[->] (id) -- (ex);
\draw[->] (ex) -- (ar);
\draw[->] (ar) -- (co);
\draw[->] (co) -- (va);
\draw[->] (va) -- (al);
\end{tikzpicture}%
}
\end{center}

\subsection{Métodos utilizados}
\begin{itemize}
    \item Entrevistas estructuradas
    \item Lecciones aprendidas
    \item Documentación técnica
    \item Videos y tutoriales
    \item Modelado de procesos (BPMN)
    \item Sistemas documentales
\end{itemize}

\subsection{Ejemplo aplicado}
En una empresa de software, el conocimiento de un desarrollador senior es documentado mediante entrevistas y diagramas, permitiendo su reutilización por otros miembros del equipo.

\section{Ciclo de Gestión del Conocimiento}

\subsection{Descripción general}
El ciclo de gestión del conocimiento describe las etapas mediante las cuales el conocimiento es generado, almacenado, compartido y utilizado dentro de una organización. Este proceso es continuo y dinámico.

\subsection{Etapas del ciclo}

\subsubsection{Captura o creación}
Se genera conocimiento a partir de la experiencia, investigación e interacción.

\subsubsection{Almacenamiento}
Se organiza en repositorios como bases de datos o sistemas documentales.

\subsubsection{Distribución}
Se comparte mediante intranets, capacitaciones y plataformas colaborativas.

\subsubsection{Uso}
Se aplica en la resolución de problemas y toma de decisiones.

\subsubsection{Retroalimentación}
Se evalúa y mejora el conocimiento.

\subsection{Diagrama del ciclo}
\begin{center}
\begin{tikzpicture}[
    node distance=2.5cm and 3cm, 
    every node/.style={
        draw, 
        circle, 
        align=center, 
        minimum size=2.2cm, 
        font=\footnotesize,
        fill=white
    },
    >=stealth, 
    thick
]

% 1. Definición única de todos los nodos
\node (c1) {Captura /\\ Creación};
\node (c2) [right=of c1] {Almacena-\\miento};
\node (c3) [below=of c2] {Distribución};
\node (c4) [below=of c1] {Retroali-\\mentación}; % Posicionado bajo c1 para cerrar el cuadro
\node (c5) [left=of c3] {Uso}; % Posicionado a la izquierda de c3 (bajo c4)

% 2. Dibujo de las flechas siguiendo el flujo lógico
\draw[->] (c1) -- (c2);
\draw[->] (c2) -- (c3);
\draw[->] (c3) -- (c5);
\draw[->] (c5) -- (c4);
\draw[->] (c4) -- (c1);

\end{tikzpicture}
\end{center}



\subsection{Relación entre etapas}
Las etapas están interconectadas formando un ciclo continuo que impulsa el aprendizaje organizacional y la mejora constante.

\subsection{Ejemplo práctico}
En soporte técnico:
\begin{enumerate}
    \item Se resuelve un problema complejo.
    \item Se documenta la solución.
    \item Se comparte con el equipo.
    \item Se reutiliza en casos similares.
    \item Se mejora con nuevas experiencias.
\end{enumerate}

\section{Conclusión}
La gestión del conocimiento es un componente esencial en la ingeniería de sistemas, ya que permite convertir experiencias individuales en activos organizacionales reutilizables. La integración del conocimiento tácito y explícito, junto con su adecuada gestión mediante ciclos estructurados, favorece la innovación, la eficiencia y la sostenibilidad organizacional.

Este enfoque es directamente aplicable al modelado de procesos, permitiendo identificar flujos de información, actores clave y puntos críticos en los que el conocimiento es generado y utilizado.

\end{document}
